\section{Conclusiones}
\label{paper:conclusiones}
La extrema y media razón holográfica queda formulada mediante una estructura generadora, una evaluación aritmética reversible y una ley de transporte con recuperación de información. El registro dodecafásico \(K\), producido por APP--TRIT--TPK, conserva su orden en la carta de fase y enlaza la clausura aritmética con la representación incidencial. Su lectura
\[
\kappa=
\frac{234543140729659824621058914794146601}{10^{36}-1}
\]
recupera exactamente las doce componentes: la multiplicación por el denominador restituye el entero y la base mil determina sus bloques. Este resultado fija el sentido matemático de la constante holográfica de acoplamiento aritmético--geométrico. Su valor es una representación exacta del registro ordenado y participa, junto con las coordenadas regionales de \(\pi,\varphi,e\), en la determinación de \(\alpha_A\).

La misma estructura permite identificar operaciones. Los doce desplazamientos cíclicos de \(K\) forman una base cuya matriz de Gram tiene extremos espectrales \(589609\) y \(6263^2\). Las respuestas a esa base determinan un operador lineal con factor de amplificación del error \(1/\sqrt{589609}\); una respuesta vectorial basta en la clase que conmuta con el desplazamiento. La representación incidencial transporta el registro a su imagen especificada con inversa explícita. Dos observaciones de memoria recuperan su parte centrada con cota óptima \(9/4\), y la carga uniforme completa el registro. Evaluación, incidencia e identificación se componen así sobre un mismo objeto, conservando en cada paso los datos requeridos por la inversa.

La coincidencia de las publicaciones entera y analítica de \(\alpha\) se establece a toda profundidad. La carta analítica prolonga el truncamiento de orden nueve mediante una recurrencia de razón \(1/729\), cuyo coeficiente inicial se obtiene de la precoordenada ya generada. La convergencia uniforme, la cota positiva de la derivada y los intervalos racionales encajados identifican la misma raíz simple y los mismos prefijos canónicos en ambas publicaciones. El teorema determina su compatibilidad completa dentro de esa construcción correlacionada.

El resultado general que sostiene la conservación es la identidad bilateral
\[
 (a+1)(aB-C)^*(aB-C)
 +a\bigl((a+1)B+C\bigr)^*\bigl((a+1)B+C\bigr)
 =\bigl((a+1)^3+a^3\bigr)I,
\]
para \(a>0\) y dos isometrías \(B,C\) con dominio y codominio comunes. La cancelación de los términos cruzados demuestra la igualdad en dimensión arbitraria. En la especialización \(a=8\), el factor ternario \(73\) conserva su interpretación como norma algebraica e índice reticular, y satisface \(10\cdot73=729+1\). El refinamiento de etiquetas enteras admisibles tiene una inversa por división euclídea, con profundidad y convención de acarreo especificadas; su levantamiento a las bases de estados es unitario y se compone con el balance.

La ley permite determinar también qué cambia al observar. Para transporte relativo unitario \(R\), el lector diagonal transforma el observable de la carta común según
\[
 \mathcal T(G)=\frac{72}{73}G+\frac1{73}R^*GR.
\]
Sus puntos fijos son exactamente los observables que conmutan con \(R\). El observable transportado sobre el estado completo conserva, en cambio, la lectura inicial por su fórmula de restitución. Las identidades exteriores prolongan este resultado a cada grado admisible, conservando los sectores puramente observados, los registrados y los mixtos. Las contribuciones de área y volumen quedan incluidas en una ley de formas, con las métricas y dimensiones correspondientes.

La recuperación se mantiene a profundidad arbitraria mediante una política explícita de archivo. En el caso nonádico, el operador terminal satisface \(\|R_N\|^2\leq(729/1241)^N\). El archivo infinito es, por tanto, una isometría y su adjunto reconstruye el estado con norma operatoria uno. Para cada profundidad positiva, el archivo finito admite una inversa exacta corregida por su Gram, con cota
\[
 \|L_N\|\leq\bigl(1-(729/1241)^N\bigr)^{-1/2}.
\]
Estas fórmulas cuantifican la estabilidad y separan la inversión exacta de la aproximación obtenida al truncar una serie de recuperación. En las dinámicas generales, la memoria y el terminal superviviente conservan conjuntamente el estado; la política indicada demuestra el régimen en que toda su norma se transfiere al archivo.

Las realizaciones físicas desarrolladas aplican estas operaciones con sus datos explícitos de acción, orientación, escala y representación. En particular, para el protocolo simpléctico declarado con dos entradas gaussianas puras, centradas e idénticas y un número finito de modos, la componente antilineal de la memoria fija el espectro simpléctico de la reducción, manteniendo todas las ocupaciones de Fock. El lazo elíptico--parabólico produce pureza \(9/11\), espectro \((9/10)10^{-j}\) y entropía adimensional \((10/9)\log10-\log9\). La conservación global y la mezcla de la observación reducida se describen dentro del mismo protocolo. Las realizaciones electromagnéticas, gravitatorias, térmicas y variacionales mantienen los operadores y normalizaciones de sus respectivas pruebas.

La construcción común del continuo conserva las cinco operaciones correlativas y la compatibilidad de sus refinamientos. La clausura de familias de historias de la sección~\ref{sec:cierre-cardinal-nativo} establece, en el universo genealógico allí definido, la igualdad cardinal interna de su recta con \(\aleph_1\) de ese universo. La afirmación conserva el dominio y las operaciones de pertenencia que intervienen en sus dos inyecciones y en su identificación final.

El censo del emisor residual hace explícita la base finita de una parte de esa cadena: reúne \(44\) firmas, \(18\) aditivas y \(26\) multiplicativas, con sus multiplicidades y la regla de reconstrucción de las fibras. Las \(648\) pertenencias de condiciones iniciales, \(324\) en cada hoja, se conservan en los datos del suplemento. Los siete teoremas comprobados en Lean certifican las identidades aritméticas y sus consecuencias finitas declaradas; las demostraciones de operadores, límites e incidencias se desarrollan en el texto. Censo, cálculo y formalización fijan así sus respectivos objetos de comprobación.

El resultado central es una caracterización calculable de la conservación evolutiva de la unidad: la estructura determina una evaluación, sus transformaciones distribuyen la información y las inversas precisan cómo recuperarla. La constante de acoplamiento, las leyes de observables y las cotas de reconstrucción expresan esa continuidad mediante ecuaciones y operaciones, desde el registro aritmético hasta las realizaciones especificadas.
